%-------------------------
% Resume in LaTeX -- "Jake's Resume"
% Author  : Jake Gutierrez
% Source  : https://github.com/jakegut/resume (fetched verbatim, then adapted)
% Based off of: https://github.com/sb2nov/resume
% License : MIT -- Copyright (c) 2020 Jake Gutierrez
%
%   Permission is hereby granted, free of charge, to any person obtaining a copy
%   of this software and associated documentation files (the "Software"), to deal
%   in the Software without restriction, including without limitation the rights
%   to use, copy, modify, merge, publish, distribute, sublicense, and/or sell
%   copies of the Software, and to permit persons to whom the Software is
%   furnished to do so, subject to the following conditions:
%
%   The above copyright notice and this permission notice shall be included in
%   all copies or substantial portions of the Software.
%
%   THE SOFTWARE IS PROVIDED "AS IS", WITHOUT WARRANTY OF ANY KIND, EXPRESS OR
%   IMPLIED. See .vendor/jake_resume_LICENSE for the full text.
%
%-------------------------
% ADAPTATIONS FOR THIS PROJECT
%
%   1. Content replaced with Jinja2 markers using the LaTeX-safe delimiters from
%      spec 6.1. See latex/env.py for the delimiter table. (This banner cannot
%      spell the markers out: Jinja parses its own delimiters even inside a
%      LaTeX comment, because it never sees the file as LaTeX.)
%   2. Every interpolated value passes through the `tex` filter. This is checked
%      mechanically by tests/test_env.py -- an unfiltered one fails the suite.
%   3. Upstream's "%%%%%% RESUME STARTS HERE" banner was rewritten: `%%` is
%      Jinja's line_statement_prefix, so that line would be parsed as code.
%   4. NOTHING about the page geometry is exposed as a variable. Margins, font
%      size and \vspace are fixed here on purpose (CLAUDE.md rule 3): the agent
%      controls how much content there is, never how tightly it is crammed.
%-------------------------

\documentclass[letterpaper,11pt]{article}

\usepackage{latexsym}
\usepackage[empty]{fullpage}
\usepackage{titlesec}
\usepackage{marvosym}
\usepackage[usenames,dvipsnames]{color}
\usepackage{verbatim}
\usepackage{enumitem}
\usepackage[hidelinks]{hyperref}
\usepackage{fancyhdr}
\usepackage[english]{babel}
\usepackage{tabularx}

%----------ATS TEXT EXTRACTION----------
% CLAUDE.md rule 4 / spec 6.3: these lines are what make the PDF's glyphs
% extract as real text instead of garbage. They are the highest-leverage lines
% in the file and they stay.
%
% They are also pdfTeX primitives, and tectonic's engine is XeTeX. Unguarded,
% \input{glyphtounicode} aborts the run with "Undefined control sequence
% \pdfglyphtounicode" -- tectonic ships no pdfTeX format at all, so there is no
% flag that makes the bare form work.
%
% XeTeX does not need them: it emits Unicode-mapped PDFs natively, because
% glyphtounicode exists to patch pdfTeX's 8-bit font encodings, which is a
% problem XeTeX does not have. So the guard runs them wherever they exist
% (pdflatex, reached via the latexmk and Docker rungs of the fallback chain)
% and skips them where they are both unavailable and unnecessary.
%
% The guard means the mechanism is now engine-dependent, so the *property* is
% verified directly on the built PDF instead of being assumed from these lines:
% tests/test_render_compile.py::test_pdf_extracts_as_real_text asserts that
% "AT&T", "C#", "100% uptime" and "user_id" come back out of the finished PDF.
\ifdefined\pdfgentounicode
  \input{glyphtounicode}
  \pdfgentounicode=1
\fi

\pagestyle{fancy}
\fancyhf{} % clear all header and footer fields
\fancyfoot{}
\renewcommand{\headrulewidth}{0pt}
\renewcommand{\footrulewidth}{0pt}

% Adjust margins
\addtolength{\oddsidemargin}{-0.5in}
\addtolength{\evensidemargin}{-0.5in}
\addtolength{\textwidth}{1in}
\addtolength{\topmargin}{-.5in}
\addtolength{\textheight}{1.0in}

\urlstyle{same}

\raggedbottom
\raggedright
\setlength{\tabcolsep}{0in}

% Sections formatting
\titleformat{\section}{
  \vspace{-4pt}\scshape\raggedright\large
}{}{0em}{}[\color{black}\titlerule \vspace{-5pt}]




%-------------------------
% Custom commands
\newcommand{\resumeItem}[1]{
  \item\small{
    {#1 \vspace{-2pt}}
  }
}

\newcommand{\resumeSubheading}[4]{
  \vspace{-2pt}\item
    \begin{tabular*}{0.97\textwidth}[t]{l@{\extracolsep{\fill}}r}
      \textbf{#1} & #2 \\
      \textit{\small#3} & \textit{\small #4} \\
    \end{tabular*}\vspace{-7pt}
}

\newcommand{\resumeSubSubheading}[2]{
    \item
    \begin{tabular*}{0.97\textwidth}{l@{\extracolsep{\fill}}r}
      \textit{\small#1} & \textit{\small #2} \\
    \end{tabular*}\vspace{-7pt}
}

\newcommand{\resumeProjectHeading}[2]{
    \item
    \begin{tabular*}{0.97\textwidth}{l@{\extracolsep{\fill}}r}
      \small#1 & #2 \\
    \end{tabular*}\vspace{-7pt}
}

% Like \resumeProjectHeading, but the left column wraps.
%
% Not an embellishment: a paper title is long by nature ("Transfer Learning for
% Fine-Grained Flower Classification on Low-Power Devices" plus a journal name
% runs ~110pt past the text block), and `tabular*` does not wrap -- it just
% overflows into the margin and reports an overfull hbox. `tabularx`'s X column
% is the same table with a width the cell can break inside.
\newcommand{\resumeWrappingHeading}[2]{
    \item
    \begin{tabularx}{0.97\textwidth}{@{}X@{\extracolsep{\fill}}r@{}}
      \small#1 & #2 \\
    \end{tabularx}\vspace{-7pt}
}

\newcommand{\resumeSubItem}[1]{\resumeItem{#1}\vspace{-4pt}}

\renewcommand\labelitemii{$\vcenter{\hbox{\tiny$\bullet$}}$}

\newcommand{\resumeSubHeadingListStart}{\begin{itemize}[leftmargin=0.15in, label={}]}
\newcommand{\resumeSubHeadingListEnd}{\end{itemize}}
\newcommand{\resumeItemListStart}{\begin{itemize}}
\newcommand{\resumeItemListEnd}{\end{itemize}\vspace{-5pt}}

%-------------------------------------------
% RESUME STARTS HERE
%-------------------------------------------

\begin{document}

%----------HEADING----------
\begin{center}
    \textbf{\Huge \scshape John Doe} \\ \vspace{1pt}
    \small (415) 555-0148 $|$ Oakland, CA $|$
    \href{mailto:john.doe@example.com}{\underline{john.doe@example.com}}
    $|$ \href{https://linkedin.com/in/johndoe}{\underline{linkedin.com/in/johndoe}}
    $|$ \href{https://github.com/johndoe}{\underline{github.com/johndoe}}
\end{center}

% Sections are emitted in the order latex/context.py computed for this profile's
% career stage -- education first while a degree is the strongest thing on the
% page, work first once it is not. Nothing here hardcodes an order, and a section
% whose entries all lost the line budget renders nothing at all.
%
% Written as LaTeX comments rather than Jinja line comments on purpose: a Jinja
% line comment removes its text but leaves the newline behind, and five stray
% blank lines here are five paragraph breaks LaTeX will happily act on. (As with
% the banner above, this note cannot spell that prefix out -- Jinja would eat the
% rest of the line, which is exactly the hazard being described.)
%----------EXPERIENCE----------
\section{Experience}
  \resumeSubHeadingListStart
    \resumeSubheading
      {Senior Backend Engineer}{Feb. 2023 -- Present}
      {Halvorsen \& Bright R\&D}{Oakland, CA}
      \resumeItemListStart
        \resumeItem{Cut p95 checkout latency from 820ms to \textasciitilde{}50ms with a Redis read-through cache and by removing N+1 queries across 12 endpoints}
        \resumeItem{Held 100\% uptime on the payments ingest for 14 months using replay-safe Kafka consumers and a dead-letter queue}
        \resumeItem{Repartitioned the events table on user\_id, taking the slowest tenant report from 41s to \textless{}100ms at \textgreater{}2M rows/day}
        \resumeItem{Migrated 9 services from EC2 to Kubernetes on EKS, cutting median deploy time from 25 to 4 minutes}
      \resumeItemListEnd
    \resumeSubheading
      {Backend Engineer}{Aug. 2020 -- Jan. 2023}
      {Northwind Data Systems}{Remote}
      \resumeItemListStart
        \resumeItem{Built a C\#/.NET reconciliation service for an AT\&T billing contract that cleared \$1.2M of disputed invoices in one quarter}
        \resumeItem{Cut monthly AWS spend 38\% by right-sizing RDS and moving cold events to S3 under a Terraform lifecycle policy}
        \resumeItem{Cut CI feedback from 22 to 6 minutes by sharding pytest across GitHub Actions runners with a stdout \textbar{} jq failure summary}
      \resumeItemListEnd
  \resumeSubHeadingListEnd
%----------PROJECTS----------
\section{Projects}
    \resumeSubHeadingListStart
      \resumeWrappingHeading
          {\textbf{LedgerLint} $|$ \emph{Python, DuckDB, pandas, Testing}}{Nov. 2023 -- May 2024}
          \resumeItemListStart
            \resumeItem{Built a DuckDB-backed linter that flags unbalanced double-entry journals, catching 340 defects across 6 public ledgers}
            \resumeItem{Cut a 12-minute full-ledger scan to 43s by pushing predicate filters into DuckDB instead of pandas}
          \resumeItemListEnd
      \resumeWrappingHeading
          {\textbf{TileCache} $|$ \emph{TypeScript, Node.js, Redis, Docker}}{June 2022 -- Mar. 2023}
          \resumeItemListStart
            \resumeItem{Shipped a Node.js \textasciicircum{}18 tile cache serving 4.1M requests/day at a 94\% hit rate and p99 under 30ms}
            \resumeItem{Added a Redis invalidation path keyed on \{z\}/\{x\}/\{y\} so stale tiles expire within 2 seconds of a source update}
          \resumeItemListEnd
      \resumeWrappingHeading
          {\textbf{VitalsBoard} $|$ \emph{React, TypeScript, FastAPI, PostgreSQL, Observability}}{May 2021 -- Dec. 2021}
          \resumeItemListStart
            \resumeItem{Built a React and FastAPI dashboard rendering 18-month vitals trends for \textasciitilde{}3k synthetic patients from a single query}
            \resumeItem{Instrumented it with OpenTelemetry traces, cutting median debug time on data gaps from 40 to 8 minutes}
          \resumeItemListEnd
    \resumeSubHeadingListEnd
%----------TECHNICAL SKILLS----------
\section{Technical Skills}
 \begin{itemize}[leftmargin=0.15in, label={}]
    \small{\item{
     \textbf{Languages}{: Python, C\#, TypeScript, SQL, Go} \\
     \textbf{Frameworks}{: FastAPI, Node.js, React, .NET} \\
     \textbf{Data}{: PostgreSQL, Redis, Kafka, DuckDB, pandas} \\
     \textbf{Infrastructure}{: AWS, Docker, Kubernetes, Terraform, GitHub Actions} \\
     \textbf{Practices}{: Caching, Query Optimization, Performance Profiling, Observability, API Design, CI/CD, Testing, Data Modeling, Cost Optimization, Incident Response}
    }}
 \end{itemize}
%----------EDUCATION----------
\section{Education}
  \resumeSubHeadingListStart
    \resumeSubheading
      {University of California, Santa Cruz}{Santa Cruz, CA}
      {B.S. Computer Science}{Sept. 2016 -- June 2020}
  \resumeSubHeadingListEnd

%-------------------------------------------
\end{document}
